As mentioned, one key feature of what our platform can provide is that we can naturally see a lot of emergent phenomena that matter for social evolution regarding morality. These phenomena were abstracted away in the traditional mathematical models. But on our platform, they will surface on their own to deepen our understanding. Those phenomena or topics were traditionally a subject of research areas on their own, but now we can study them in a unified framework. 

We list some of the observed phenomena and identify some of the theories that are related to them in Table~\ref{tab:project_findings}. This list is definitely not exhaustive. We hope this can provide a good starting point for future researchers to discover more phenomena and theories.

We also encourage researchers to use our platform as a new way to study these phenomena and theories.

\begin{table*}[tbp]
\centering
\caption{Discovered Phenomena and Related Theories}
\label{tab:project_findings}
\small
\setlength{\tabcolsep}{4pt}
% \renewcommand{\arraystretch}{0.8}  % Reduce vertical spacing between rows
\begin{tabular}{p{7cm}p{7cm}}
\hline
\textbf{Phenomena Findings from Experiments} & \textbf{Related Theories} \\
\hline
\parbox[t]{7cm}{Coordination is costly: \\ 
\hspace*{1em} $\bullet$ Communication takes time and can reduce the time for other important things.} & 
\parbox[t]{7cm}{\textit{Coordination Cost Theory} \citep{simon1991bounded} \\
"Organizations face bounded rationality where coordination costs limit optimal decision-making"} \\
\hline
\parbox[t]{7cm}{Moral judgment based on actions: \\
\hspace*{1em} $\bullet$ Agents evaluate others' morality by observing how they treat third parties. \\
\hspace*{1em} $\bullet$ Actions toward others, not just toward oneself, shape moral reputation.} & 
\parbox[t]{7cm}{\textit{Moral Judgment Theory} \citep{haidt2001emotional} \\
"People make rapid moral judgments based on observed behaviors and their emotional responses" \\
\textit{Impression Formation} \citep{asch1946forming} \\
"Observers form impressions of others' character based on their actions toward third parties"} \\
\hline

\parbox[t]{7cm}{Misunderstandings can lead to major conflicts: \\
\hspace*{1em} $\bullet$ Agents may misinterpret others' intentions or actions, leading to unnecessary conflicts. \\
\hspace*{1em} $\bullet$ Limited communication can cause agents to make incorrect assumptions about others' moral types or goals.} & 
\parbox[t]{7cm}{\textit{Communication Theory} \citep{shannon1948mathematical} \\
"Information transmission is inherently imperfect, leading to potential misunderstandings and conflicts" \\
\textit{Conflict Resolution} \citep{deutsch1973resolution} \\
"Many conflicts arise from misperceptions and misunderstandings rather than actual incompatible goals"} \\
\hline

\hline
\parbox[t]{7cm}{Predictable morality acts as a reputation/clear signal: \\
\hspace*{1em} $\bullet$ Universal moral agents, by rejecting violence completely, benefit from being clearly understood. \\
\hspace*{1em} $\bullet$ Such behavior is costly though, making them subject to exploitation} &
\parbox[t]{7cm}{\textit{Costly Signaling Theory} \citep{gintis2001costly} \\
"Consistently moral behavior, despite its costs, serves as an honest signal of cooperative intent, reducing the risk of being misjudged as a threat." \\
\textit{Indirect Reciprocity} \citep{nowak2005evolution} \\
"A clear reputation for cooperation, built through predictable actions, is essential for reciprocal altruism and protects an agent from being wrongly punished."} \\
\hline
\parbox[t]{7cm}{Universal moral agents get exploited: \\
\hspace*{1em} $\bullet$ Agents who never retaliate or punish others' bad behavior become targets of exploitation. \\
\hspace*{1em} $\bullet$ Their unconditional cooperation makes them vulnerable to free-riders.} & 
\parbox[t]{7cm}{\textit{Altruistic Punishment} \citep{fehr2002altruistic} \\
"Cooperation requires punishment of defectors; pure altruism without retaliation is vulnerable to exploitation" \\
\textit{Strong Reciprocity} \citep{bowles2004evolution} \\
"Evolutionary success requires both cooperation and punishment of non-cooperators"} \\
\hline
\parbox[t]{7cm}{Group membership is contested: \\
\hspace*{1em} $\bullet$ Agents might not agree on who is in the group that can share resources.} & 
\parbox[t]{7cm}{\textit{Social Identity Theory} \citep{tajfel1979individuals} \\
"Group boundaries are fluid and contested, with membership determined by shared identity markers and mutual recognition"} \\
\hline
\parbox[t]{7cm}{Distribution methods are complex: \\
\hspace*{1em} $\bullet$ How to distribute? Distribute evenly, based on contribution, harm taken, need, can affect both the success of the end result and each other's judgment.} & 
\parbox[t]{7cm}{\textit{Distributive Justice} \citep{konow2003which} \\
"Fairness judgments depend on multiple principles, including equality, need, and contribution"} \\
\hline
\parbox[t]{7cm}{Careful planning is important: \\
\hspace*{1em} $\bullet$ Reproduction schedule is important. Too frequent can cause both parents and children to die.} & 
\parbox[t]{7cm}{\textit{Life History Theory} \citep{kaplan1992evolution} \\
"Organisms face trade-offs between current and future reproduction, with timing being crucial for survival"} \\
\hline
\parbox[t]{7cm}{Tendency to cooperate can sometimes have negative effect: \\ 
\hspace*{1em} $\bullet$ Moral agents have a tendency to collaborate to acquire resources, but in some particular setting (with competition, resources being in some way), taking faster action instead of collaboration may be more crucial. \\
\hspace*{1em} $\bullet$ Moral agents tend to agree to collaborate to hunt, but they might not be in a good position to hunt.} & 
\parbox[t]{7cm}{\textit{Cooperation Dilemmas} \citep{bowles2004evolution} \\
"Cooperation can be maladaptive when individual action would yield higher returns"} \\
\hline
\parbox[t]{7cm}{Moral agents' mutual dependency sometimes leads to disaster end: \\
\hspace*{1em} $\bullet$ Moral agents tend to trust others to help them later, but the others may also think the same, and none have the extra capacity to help.} & 
\parbox[t]{7cm}{\textit{Trust and Cooperation} \citep{fehr2002altruistic} \\
"Altruistic punishment can maintain cooperation but may lead to cascading failures when trust is misplaced"} \\
\hline
\parbox[t]{7cm}{Mutual reinforcing / social pressure: \\
\hspace*{1em} $\bullet$ When some agents reproduce, others feel compelled to do so too, even though their HP was not very high.} & 
\parbox[t]{7cm}{\textit{Social Learning Theory} \citep{bandura1977social} \\
"Social learning and imitation can lead to behavioral contagion even when not optimal for individuals"} \\
\hline
\end{tabular}
\end{table*}